\section{Introduction}\label{sec:introduction}

The Ashkin--Teller model originates in the four-component lattice model
of Ashkin and Teller \cite{AshkinTeller1943}. Its current representation
raises a geometric question beyond convergence of the associated height:
what is the joint continuum law of all primal and dual current clusters,
and how are these collections coupled to the field? We prove the joint
limit in \cref{thm:main} for every fixed point on the critical line,
including the four-state Potts endpoint. This is the height and
current-cluster assertion formulated by Alcalde L\'opez, Heeney, and Lis
\cite[Conjecture~1.2]{AlcaldeLopezHeeneyLis2026}.

Fix $J>0$ and $U\le J$ with $\sinh(2J)=e^{-2U}$. Our zero-boundary
Gaussian free field has covariance $2\pi(-\Delta_D)^{-1}$, with
logarithmic singularity $-\log|z-w|$. In this normalization put
\[
 \lambda=\frac\pi2,\qquad a=\lambda\sqrt g,\qquad
 g=\frac8\pi\arcsin\!\left(\frac{\sqrt{1+e^{4U}}}{2}\right).
\]
Thus $4/3<g\le4$ and $\lambda<a\le2\lambda$. For the wired current,
the first limiting cluster is the symmetric two-valued local set with
labels $\pm a$. For the free current, first expose $\CLE_4$ holes and
retain an asymmetric two-valued split inside each hole; repeat in all
complementary components. Both constructions use the same field.
The theorem retains every nesting depth and the distinguished wired
boundary cluster. It includes tightness and identification of the entire
compact supports in positive-diameter Hausdorff matching topology.

\subsection*{Prior work and the scope of the result}
The finite spin, percolation, and height representation used here is due
to Lis \cite[Sections~1.1--1.3 and~2.2]{Lis2022Spins}; his current expansion
\cite[Section~2]{Lis2022Boundary} gives the precise odd/even-edge weights.
A noncrossing primal--dual contour representation appears earlier in Rys
\cite[Section~3]{Rys1963}. We recall and verify the finite coupling,
including its boundary convention, in \cref{sec:model}.

At $U=0$, the currents are double random currents and $g=2$.
Duminil-Copin and Lis \cite[Theorems~1.1--1.2]{DuminilCopinLis2019}
identified the finite signed nesting field with a dimer height.
Duminil-Copin, Lis, and Qian \cite{DRC2025} established conformally
invariant limits for the nesting field and recursive cluster-boundary
families. Their introductory statements cover Jordan domains, while the
detailed joint boundary-family Theorems~6.2 and~6.4 assume a $C^1$
boundary and use injective matching of loops. Alcalde L\'opez, Heeney,
and Lis \cite[Theorem~1.14 and Appendix~B]{AlcaldeLopezHeeneyLis2026}
state the joint limit of the height and both complete compact
current-cluster collections in Jordan domains, using the noninjective
positive-diameter matching topology considered here. Their
Conjecture~1.2 asks for the critical-line extension obtained by replacing
$\sqrt2$ with $\sqrt g$.

Alcalde L\'opez and Sep\'ulveda
\cite[Theorems~1.4 and~9.5, and Lemma~9.1]{AlcaldeLopezSepulveda2026} study a stronger
joint decomposition at the Ising point involving the XOR field, cluster
measures, and signs, with the Gaussian field identified with the nesting
height. Their introductory theorem is stated for Jordan domains; their
detailed joint convergence theorem adds a boundary Minkowski-dimension
hypothesis. Those additional field and measure limits are
outside our statement. Our conclusion concerns the height and complete
current supports for every fixed $4/3<g\le4$.

The analytic input is related to the six-vertex model with weights
$(1,1,c)$, where $c=\sqrt{1+e^{4U}}\in(1,2]$.
Duminil-Copin, Kozlowski, Lammers, and Manolescu \cite{DKLM2026}
prove a height scaling limit for $\sqrt3\le c\le2$ and establish several
inputs on the larger interval $1\le c\le2$: logarithmic bounds for even
increment moments, one-parity spin-percolation circuits, a specified
flat-exhaustion plane law, and a free-energy identification of the
normalization conditional on a Gaussian plane limit. Odd plane correlations
vanish by height reversal. We use these larger-range results. The covariance, screening,
and Gaussianity arguments needed for the whole interval are given here;
the free-energy formula is applied only after its convergence premise
has been established. A related full-plane height theorem for the balanced
six-vertex state is proved in \cite[Theorem~1.1]{BalancedSixVertex2026}. Its unit
height corresponds to $K=2H$ here, and its multiplier
$1/\arcsin(c/2)$ agrees with ours after this normalization. This
comparison does not replace the local Gaussianity proof or the
random-domain and cluster identifications below.

For the reflection argument, coordinate and diagonal tests for
six-vertex weights already appear in
\cite[Section~6]{FrohlichIsraelLiebSimon1980}. We use the general
reflection-positive framework of \cite{FrohlichIsraelLiebSimon1978};
for the reflected Hilbert-space construction and positive-transfer
representation, see \cite[Sections~3.1--3.3 and~4.1]{Usui2012}.
The height transfer in \cite[Section~4.5.3]{DKLM2026} is a closer
model-specific precedent and permits negative one-step spectrum.
The transfer, screening, and collision strategy also has a related
implementation for integer Lipschitz heights with two-arc boundary data
in \cite[Sections~5, 7, and~9--10]{Triangular}. We prove the square-frame
and random-chamber arguments required here.

The continuum identification uses thin two-valued local sets
\cite{ASW2017,AS2018}, first-passage sets \cite{ALS2017,ALS2018}, and
the canonical excursion decomposition \cite{ALS2023}. The recursive
exploration and thinness architecture also follows the double-current
work \cite{DRC2025}. In the present coupling, we must additionally
identify which discrete odd holes belong to a single current cluster,
exclude macroscopic pieces invisible to those holes, and compare passage
tests after a second coupled random exploration. We prove these steps
with the conditioning and approximation hypotheses stated in the
corresponding propositions. A proof strategy follows the formal theorem
in \cref{subsec:proof-strategy}.
